% Applied Software Development Lab -- Lab Handout 5: The word class makes an appearance
% Copyright (c) Michael Lehn. Licensed under CC BY-SA 4.0.
%
% Region coordinates (bp from the top of the page) follow the page layout of
% the original German handout "Laborblatt 5".
\documentclass{aslab}
\sethandout{5}{The word class makes an appearance}
\renewcommand*\handouttitle{The word \codearg{class} makes an appearance}
\graphicspath{{../figures/handout05/}{../figures/shared/}}

%% Lists used in this handout
% "A)" items (bold titles, text at +18bp)
\newlist{letters}{enumerate}{3}
\setlist[letters]{label={\Alph*)},leftmargin=18bp,labelsep=*,align=left,
  labelindent=0pt}
% "(i)" items
\newlist{romans}{enumerate}{3}
\setlist[romans]{label={(\roman*)},leftmargin=18bp,labelsep=*,align=left,
  labelindent=0pt}
% "(I)" items
\newlist{Romans}{enumerate}{3}
\setlist[Romans]{label={(\Roman*)},leftmargin=18bp,labelwidth=15bp,
  labelsep=3bp,align=left,labelindent=0pt}
% bold "(1)" items (input lines)
\newlist{inputs}{enumerate}{3}
\setlist[inputs]{label={\bfseries(\arabic*)},leftmargin=18bp,labelsep=*,
  align=left,labelindent=0pt}
% numbered lists with a regular-weight number
\newlist{ptasks}{enumerate}{3}
\setlist[ptasks]{label={\arabic*.},leftmargin=18bp,labelsep=*,align=left,
  labelindent=0pt}
% "•" items with the text at +9bp
\newlist{dotlist}{itemize}{3}
\setlist[dotlist]{label={\textbullet},leftmargin=9bp,labelsep=*,align=left,
  labelindent=0pt}

%% Python listings with the colors of the original
%%   @..@ orange (line numbers, keywords), ~..~ cyan (built-ins, names),
%%   ^..^ red (strings, numbers), !..! magenta (import, escapes)
\lstdefinestyle{py}{style=aslab,
  moredelim=[is][\color{codeorange}]{@}{@},
  moredelim=[is][\color{codecyan}]{~}{~},
  moredelim=[is][\color{codered}]{^}{^},
  moredelim=[is][\color{codemagenta}]{!}{!}}
% typographic quotes inside a listing (as in the original code): |...|
\lstdefinestyle{curly}{style=aslab,escapechar=|}
\newcommand*\codeix{\fontsize{9bp}{12bp}\ttfamily\lineskiplimit=-\maxdimen}
\newcommand*\codex[1]{\fontsize{10bp}{#1bp}\ttfamily\lineskiplimit=-\maxdimen}

% keep the baseline pitch fixed even next to tall inline code boxes
\lineskiplimit=-\maxdimen

\begin{document}
%% ===================================================================== 1 ==
\maketitleblock

\begin{region}{124.1}
\heading{Learning objectives}
\end{region}

\begin{region}[pitch=13.3]{138.6}
{\bfseries\color{aslabblue}
\begin{dashlist}
\item Reading the serial interface with Python
\item Using control structures ({\mdseries\code{if}}, {\mdseries\code{for}},
  {\mdseries\code{while}}) and simple data types
\item A step-by-step introduction to object-oriented programming (including
  the motivation, but without overkill -- the rest comes later)
\end{dashlist}\par}
\end{region}

\begin{region}{204.6}
\heading{Step 1: Cloning your own repository and building the game controller}
\end{region}

\begin{region}[pitch=14.2]{233.2}
\textbf{Goal:} Picking up where we left off last time -- starting with the
microcontroller code.
\end{region}

\begin{region}[pitch=13.2]{275.0}
On the breadboard, rebuild the game controller from Lab Handout~4, connect it
to the microcontroller, and connect the microcontroller to your computer via
USB.
\end{region}

\begin{region}[pitch=13.2]{314.5}
As always:
\begin{dashlist}
\item Open two terminals (left for editing, right for running commands).
\item In addition, open a web browser and go to your Git repository at\\
  \code{https://gitlab.uni-ulm.de/USERNAME/asp-ulm}
\end{dashlist}
\end{region}

\begin{region}{393.3}
\subheading{Left terminal (editing)}
\end{region}

\begin{region}[pitch=13.35]{407.6}
\begin{letters}
\item \textbf{Clone the repository and change the directory}\\
  Clone the repository with:\\
  \code{git clone https://gitlab.uni-ulm.de/USERNAME/asp-ulm}\\
  Change the directory:\\
  \code{cd asp-ulm/controller}
\end{letters}
\end{region}

\begin{region}[pitch=13.6]{474.9}
\begin{letters}[start=2]
\item \textbf{Open the editor}\\
  Open the file \code{controller.ino} in the editor of your choice.
\end{letters}
\end{region}

\begin{region}[pitch=13.35]{502.7}
\begin{letters}[start=3]
\item \textbf{Think ;-)}\\
  Read the code carefully and think about what should happen when you
  press the push button or turn the potentiometer.
\end{letters}
\end{region}

\begin{region}{568.8}
\subheading{Right terminal (running commands)}
\end{region}

\begin{region}[pitch=12.5]{581.8}
\begin{letters}
\item \textbf{Change the directory}\\
  \code{cd asp-ulm/controller}
\item \textbf{Make the USB port visible to WSL}\\
  \code{usb_check}
\item \textbf{Flash the microcontroller}\\
  \code{make upload}
\item \textbf{Start the serial monitor}\\
  \code{make monitor}
\item \textbf{Test}\\
  Check whether the expected behavior occurs when you press the push button
  and turn the potentiometer.
\end{letters}
\end{region}

\nextpage
%% ===================================================================== 2 ==
\begin{region}{69.4}
\heading{Step 2: A serial monitor of our own making}
\end{region}

\begin{region}{98.0}
\textbf{Goal:} Reading the serial interface with a Python program.
\end{region}

\begin{region}[pitch=13.5]{124.5}
What is special about our own version: the output is not produced
character by character, but only once a complete line (that is, up to the
line break) has been received. As soon as that is the case, \textbf{this
line is output all at once}. As long as that is not the case, nothing
happens.
\end{region}

\begin{region}{205.5}
\subheading{Preparation}
\end{region}

\begin{region}[pitch=13.55]{233.0}
\begin{tasks}
\item \textbf{Left terminal}
  \begin{letters}
  \item \textbf{Change the directory}\\
    \code{cd ../python}
  \item \textbf{Open the editor}\\
    Open the file \code{serial-reader.py} in the editor of your choice.
  \item \textbf{Think ;-)}\\
    Read the code carefully and think about what should happen when you
    press the push button or turn the potentiometer. If necessary, go
    through Step~6 of Lab Handout~4 again to make sure you understand:
    \begin{romans}
    \item What do \code{ser.in_waiting} and \code{ser.read} do?
    \item What do the remaining statements do?
    \item How are the loops nested -- and why?
    \end{romans}
  \end{letters}
\end{tasks}
\end{region}

\begin{region}[pitch=13.65]{396.1}
\begin{tasks}[start=2]
\item \textbf{Right terminal}
  \begin{letters}[label={\bfseries\Alph*)}]
  \item \textbf{Change the directory}\\
    \code{cd ../python}
  \item \textbf{Test}\\
    Start the program and check whether the expected behavior occurs when
    you press the push button and turn the potentiometer.
  \end{letters}
\end{tasks}
\end{region}

\begin{region}[pitch=13.5]{490.9}
You should now be at the point where individual bytes are received. For
example, \code{ch} is output as \code{b'X'} -- an object of type bytes, i.e.,
a data type for ``raw'' binary data. How to convert \code{ch} into readable
text is something we will explore experimentally in a Python shell in a
moment.
\end{region}

\begin{region}{558.5}
\subheading{Tasks}
\end{region}

\begin{region}[pitch=13.3]{585.2}
\begin{tasks}
\item If you have not already done so, stop the Python program in the
  right terminal with \code{Control-C}.
\end{tasks}
\end{region}

\begin{region}{611.9}
\begin{tasks}[start=2]
\item Start the interactive Python shell with \code{python}
\end{tasks}
\end{region}

\begin{region}[pitch=13.2]{625.2}
\begin{tasks}[start=3]
\item The Python shell works much like the terminal shell you use to
  navigate the file system or run commands -- except that here you type in
  Python statements.
  \blank
  Enter the following, one after the other:
  \begin{inputs}
  \item \code{ch = b'A'}
  \item \code{print(ch)}
  \item \code{type(ch)}
  \item \code{s = bytes.decode(ch)}
  \item \code{type(s)}
  \end{inputs}
\end{tasks}
\end{region}

\nextpage
%% ===================================================================== 3 ==
\begin{region}[x=74.7,right=0]{68.5}
You should observe:
\end{region}

\begin{region}[x=74.7,right=0,pitch=14.3]{95.6}
\begin{codeblock}[basicstyle=\codex{14.3}]
>>> ch = b'A'
>>> print(ch)
b'A'
>>> type(ch)
<class 'bytes'>
>>> type(ch)
<class 'bytes'>
>>> s = bytes.decode(ch)
>>> type(s)
<class 'str'>
\end{codeblock}
\end{region}

\begin{region}[x=74.7,right=0,pitch=13.5]{251.3}
You can use \code{type(...)} to query the type of a variable. And the
function \code{bytes.decode(...)} lets you convert a bytes object such as
\code{ch} into a string.
\end{region}

\begin{region}[pitch=13.25]{279.0}
\begin{tasks}[start=4]
\item \textbf{Tip for using the Python shell:} With the \textbf{up and down
  arrow keys} you can scroll through the input history and run or edit
  commands you have already entered.
\end{tasks}
\end{region}

\begin{region}[pitch=13.2]{318.6}
\begin{tasks}[start=5]
\item Now we explore experimentally how strings can be concatenated. To do
  this, enter the following in the Python shell:
  \begin{inputs}
  \item \code{buffer = ""}
  \item \code{print(buffer)}
  \item \code{buffer = buffer + s}
  \item \code{print(buffer)}
  \item \code{buffer += s}
  \item \code{print(buffer)}
  \end{inputs}
\end{tasks}
\end{region}

\begin{region}[x=74.7,right=0,pitch=13.4]{437.4}
You see:
\begin{dotlist}
\item With \code{+} you can concatenate two strings.
\item The notation \code{buffer += s} is equivalent to
  \code{buffer = buffer + s}, but shorter and more convenient.
\end{dotlist}
\end{region}

\begin{region}[pitch=13.2]{490.7}
\begin{tasks}[start=6]
\item Now we use a for loop to simulate that we receive individual bytes one
  after the other, convert each of them into text, and append it to a
  string.
  \blank
  Enter the following:
\end{tasks}
\end{region}

\begin{region}[x=74.7,right=0,pitch=14.3]{557.3}
\begin{codeblock}[basicstyle=\codex{14.3}]
buffer = ""
for ch in [b'A', b'\n', b'X', b' ', b'1', b'\n']:
    ch = bytes.decode(ch)
    buffer += ch
print(buffer)
\end{codeblock}
\end{region}

\begin{region}[x=74.7,right=0]{642.3}
This is not quite what we will need later -- but it is heading in the right
direction.
\end{region}

\begin{region}[x=74.7,right=0,pitch=13.5]{668.7}
Here is what happens:
\begin{dotlist}
\item \code{buffer} is initially an empty string.
\item In each loop iteration, a byte (\code{ch}) is converted into a string.
\item This string is appended to \code{buffer}.
\item At the end, \code{buffer} is output once.
\end{dotlist}
Note: Although there is only one \code{print} statement, two lines are
displayed.\\
This is because \code{buffer} contains two line breaks
(\code{\n}) -- and \code{print(buffer)} outputs the entire text including
the line breaks it contains.
\end{region}

\nextpage
%% ===================================================================== 4 ==
\begin{region}[pitch=13.35]{68.8}
\begin{tasks}[start=7]
\item Now the line break (\code{\n}) should no longer be appended to
  \code{buffer}. Instead, in this case the current content of \code{buffer}
  should be output, and \code{buffer} should then be reset to an empty
  string. To do this, change the code as follows:
\end{tasks}
\end{region}

\begin{region}[x=74.7,right=0,pitch=12.1]{112.6}
\begin{codeblock}[style=curly,basicstyle=\codex{12.1}]
buffer = ""
for ch in [b'A', b'\n', b'X', b' ', b'1', b'\n']:
    ch = bytes.decode(ch)
    if ch == '\n':
        print(f"Zeile> |\textquoteleft|{buffer}|\textquoteright|")
        buffer = ""
    else:
        buffer += ch
\end{codeblock}
\end{region}

\begin{region}[x=74.7,right=0,pitch=13.5]{214.5}
The statement \code{print(f"Zeile> {buffer}")} produces formatted output:
the placeholder \code{{buffer}} is automatically replaced by the current content
of the variable \code{buffer}.
\end{region}

\begin{region}[x=74.7,right=0,pitch=12.1]{245.2}
\begin{codeblock}[style=curly,basicstyle=\codex{12.1}]
Zeile> |\textquoteleft|A|\textquoteright|
Zeile> |\textquoteleft|X 1|\textquoteright|
\end{codeblock}
\end{region}

\begin{region}[pitch=13.45]{274.3}
\begin{tasks}[start=8]
\item One more little thing: the microcontroller does not send line breaks as
  a single \code{b'\n'}, as is usual in the Unix world, but in a format that
  the Windows world also understands: a \code{b'\r'} (carriage return)
  followed by \code{b'\n'} (line feed).
\end{tasks}
\end{region}

\begin{region}[x=74.7,right=0,pitch=13.4]{319.1}
This is a perfect opportunity, though, to introduce an \code{elif} and the
keyword \code{pass} -- so that we simply ignore the character \code{\r}, as
it should be. To do this, change the code as follows:
\end{region}

\begin{region}[x=74.7,right=0,pitch=12.1]{362.9}
\begin{codeblock}[basicstyle=\codex{12.1}]
buffer = ""
for ch in [b'A', b'\r', b'\n', b'X', b' ', b'1', b'\r', b'\n']:
    ch = bytes.decode(ch)
    if ch == '\n':
        print(f"Zeile> {buffer}")
        buffer = ""
    elif ch == '\r':
        pass
    else:
        buffer += ch
\end{codeblock}
\end{region}

\begin{region}[x=74.7,right=0,pitch=13.5]{489.1}
The \code{elif} condition detects the carriage return character
(\code{'\r'}) and then skips it with \code{pass}. Everything else is
appended to \code{buffer} as before.\\
\textbf{Additional task:} Test what happens without the \code{elif} fix.
\end{region}

\begin{region}[x=74.7,right=0,pitch=13.25]{543.5}
\code{pass} in Python simply means: do nothing. This is comparable to a
\code{nop} instruction (``no operation'') in assembly language -- a
placeholder that explicitly does nothing but has to be there syntactically.
\end{region}

\begin{region}[pitch=14.2]{584.1}
\begin{tasks}[start=9]
\item \textbf{Now comes the grand finale}\\
  Now it is your turn to put the \textbf{knowledge from your experiments}
  into practice!
\end{tasks}
\end{region}

\begin{region}[x=74.7,right=0,pitch=13.4]{626.8}
\begin{dotlist}
\item \textbf{Exit the Python shell} with \code{Control-D}.
\item Modify \code{serial-reader.py} so that the received bytes are output
  line by line as text. As soon as a complete line has been received (ending
  with \code{\n}), it should be output in full. The carriage return
  character (\code{\r}) is ignored -- as before.
\item \textbf{Test your program} in the right terminal.
\item If everything works:\\
  \code{git commit -a -m "serial reader with line buffer"}\\
  \code{git push}\\
  With this, you have built your own serial non-blocking buffered line
  reader -- congratulations!
\end{dotlist}
\end{region}

\nextpage
%% ===================================================================== 5 ==
\begin{region}{82.6}
\heading{Goal of the following steps: Code organization is everything!}
\end{region}

\begin{region}[pitch=13.17]{110.2}
In the next steps, the existing code is reorganized so that it becomes
modular, reusable, and easy to extend -- for example, as a building block
for later projects (such as our game). Along the way, we introduce the
concept of object-oriented programming (OOP) step by step -- each time with a
concrete example.\\
In the end, two central functions (or methods) should be available:
\end{region}

\begin{region}[pitch=13.44]{188.9}
\begin{dashlist}
\item \code{connect(port)}\\
  Establishes the connection to the serial interface and returns an open
  Serial object, which is stored internally.
\item \code{readline()}\\
  Checks whether a complete line has been received. Returns the line as a
  string if there is one, otherwise \code{None}. The received bytes are
  buffered internally and only output as a line once a line break
  (\code{\n}) has been received. The carriage return character (\code{\r})
  is ignored.
\end{dashlist}
\end{region}

\begin{region}{309.4}
The restructuring takes place in several stages:
\end{region}

\begin{region}[pitch=13.55]{336.7}
\begin{ptasks}
\item \textbf{Modularization with global variables:}\\
  The logic is moved into two functions that use the global variables
  \code{ser} and \code{buffer}. This results in a first, easy-to-understand
  modular structure.
\item \textbf{Conversion into a class with static attributes and methods:}\\
  In the second step, the existing structure is turned into a class.
  \code{ser} and \code{buffer} become static attributes, \code{connect} and
  \code{readline} become static methods. This does not change how the code
  is used, but it brings order into the program.
\item \textbf{(Optional) Instance-based design for multiple use:}\\
  In a further step, the class could be extended so that several serial
  connections are possible at the same time -- using instance attributes
  and methods.
\end{ptasks}
\end{region}

\begin{region}[pitch=13.75]{485.9}
\subheading{No need to fear technical terms:}
Do not let words such as \emph{static attributes}, \emph{instance
attributes}, \emph{methods}, \emph{modules}, or \emph{classes} put you off.
We develop all of this step by step using concrete examples -- and you will
see: behind these ``fancy'' expressions are simple, easy-to-understand
concepts.
\end{region}

\nextpage
%% ===================================================================== 6 ==
\begin{region}{69.4}
\heading{Step 3: Functions with shared variables}
\end{region}

\begin{region}[pitch=13.2]{97.0}
Fortunately, we already have a working program that reads the serial
interface line by line. This already gives us a good idea of what the two
planned functions should look like:
\end{region}

\begin{region}[pitch=13.54]{149.4}
\begin{dashlist}
\item The function \code{connect} uses \code{ser = serial.Serial(...)} to
  establish a connection to the USB port.
\item The variable \code{ser} is then used in the function \code{readline}
  to
  \begin{dashlist}
  \item query whether new bytes are available
    (\code{ser.in_waiting > 0}), and
  \item read a byte if needed (\code{ser.read(1)}).
  \end{dashlist}
\item The function \code{readline} also uses a variable \code{buffer} to
  which new bytes are appended (after decoding). This buffer should have
  been initialized as an empty string before the first call of
  \code{readline} -- this can conveniently be done in \code{connect}.
\end{dashlist}
\end{region}

\begin{region}[pitch=14.2]{270.9}
So we need two variables -- \code{ser} and \code{buffer} -- that can be used
by both functions. The obvious solution: \textbf{global variables}.
\end{region}

\begin{region}[pitch=13.35]{312.5}
However, Python has a nasty little catch when you want to \textbf{modify}
global variables. Let us look at this using a simple proof of concept:\\
In the left terminal, create the file \code{controller.py} in an editor of
your choice, with this content:
\end{region}

\begin{region}[pitch=12]{362.5}
\begin{codeblock}[style=py,basicstyle=\codeix]
@  1 @ser = ~None~
@  2 @buffer = ~None~
@  3@
@  4 def@ ~connect~(port):
@  5@
@  6 @    ser = f^"connect to {port}"^
@  7 @    buffer = ^"empty"^
@  8@
@  9 def@ ~readline~():
@ 10@
@ 11 @    ~print~(f^"reading from {ser}"^)
@ 12 @    buffer += ^", some data"^
@ 13 @    ~print~(f^"buffer: {buffer}"^)
@ 14@
@ 15 @connect(^"/dev/ttyUSB0"^)
@ 16 @readline()
\end{codeblock}
\end{region}

\begin{region}[pitch=13.63]{569.5}
By ``proof of concept'' we mean: we use the simplest possible example to
test whether two functions can share the same variables. The goal is
\textbf{not yet} the serial interface -- for now, the point is only that
\code{connect} ``prepares'' something, and \code{readline} accesses it and
carries on.
\end{region}

\begin{region}[pitch=14.2]{624.6}
However, this does \textbf{not} work yet. We will \textbf{analyze and fix}
this step by step in a moment, and then gradually build in the desired
functionality.
\end{region}

\begin{region}{666.3}
\subheading{Tasks}
\end{region}

\begin{region}[pitch=13.3]{693.0}
\begin{ptasks}
\item In the right terminal, run the program with
  \code{python controller.py} and confirm that there is an error.
\end{ptasks}
\end{region}

\begin{region}[pitch=14.3]{720.4}
\begin{ptasks}[start=2]
\item The error message says (in a more or less understandable way) that
  in \textbf{line 12} the variable \code{buffer} is accessed, but that it
  \textbf{does not exist}. This is actually surprising -- after all, we did
  create \code{buffer} in \textbf{line 2}, outside the function!\\
  But: \textbf{by default, Python assumes} that only \textbf{local
  variables are modified} inside a function. So if you write
  \code{buffer += ...}, Python interprets
\end{ptasks}
\end{region}

\nextpage
%% ===================================================================== 7 ==
\begin{region}[x=74.7,right=0,pitch=13.7]{68.8}
this as an attempt to modify a local variable \code{buffer} --
and that one has not been defined before.\\
So we have to \textbf{tell Python explicitly} that the global variable
\code{buffer} is meant. This is done by changing \textbf{line 10} to:
\end{region}

\begin{region}[x=74.7,right=0,pitch=13.65]{123.7}
\makebox[27.5bp][l]{\color{codeorange}10}{\color{codeorange}global}
buffer\\
Run the program again now. Confirm that there is still an error message
-- \textbf{but a different one}.
\end{region}

\begin{region}[pitch=14.02]{165.3}
\begin{ptasks}[start=3]
\item \textbf{The new error message} says that the expression
  \code{buffer += ...} in \textbf{line 12} is not allowed -- because
  \code{buffer} has the value \code{None} at this point. When
  \code{connect} is called, a variable \code{buffer} is indeed set in \textbf{line 7},
  \textbf{but not the global variable} that we created in \textbf{line 2}.
  Instead, Python simply created a \textbf{local variable} \code{buffer} in
  the function \code{connect} -- just as it always does by default. The
  global variable \code{buffer} therefore remains \textbf{unchanged} at
  \code{None}.
\end{ptasks}
\end{region}

\begin{region}[x=74.7,right=0,pitch=17.7]{249.4}
\textbf{Solution:} In the function \code{connect}, too, we have to say
explicitly that we want to modify the global variable \code{buffer}. And by
now you know how to do that \emoji{1F609}
\end{region}

\begin{region}[pitch=14.2]{297.0}
\begin{ptasks}[start=4]
\item \textbf{There is no longer an error message}, but the program does
  \textbf{not do what it is supposed to do}. The output is:
\end{ptasks}
\end{region}

\begin{region}[x=74.7,right=0,pitch=12.1]{337.2}
\begin{codeblock}[basicstyle=\codex{12.1}]
reading from None
buffer: empty, some data
\end{codeblock}
\end{region}

\begin{region}[x=74.7,right=0]{376.0}
\textbf{What we would expect, however:}
\end{region}

\begin{region}[x=74.7,right=0,pitch=12.1]{402.0}
\begin{codeblock}[basicstyle=\codex{12.1}]
reading from connect to /dev/ttyUSB0
buffer: empty, some data
\end{codeblock}
\end{region}

\begin{region}[x=74.7,right=0,pitch=13.85]{440.1}
In the function \code{connect}, we did correctly modify the global
\code{buffer} -- but \code{ser} was also only initialized locally. As a
result, the global variable \code{ser} remains \code{None}, and that is
exactly what we see in the output of \code{readline}.\\
\textbf{Solution:} The variable \code{ser}, too, must be explicitly declared
as \textbf{global} in the function \code{connect}. This is done by writing
\textbf{line 5} as follows:
\end{region}

\begin{region}[x=74.7,right=0]{523.6}
\begin{codeblock}[style=py,basicstyle=\codex{12}]
@ 5 @    @global@ buffer, ser
\end{codeblock}
\end{region}

\begin{region}[x=74.7,right=0,pitch=14.3]{551.6}
\textbf{The nastiest part of it all:} In the function \code{readline}, we
also access \code{ser} in \textbf{line 11} -- without writing global
\code{ser} there. \textbf{And yet it works!}
\end{region}

\begin{region}[x=74.7,right=0,pitch=13.7]{593.4}
\textbf{Why?} Because in \code{readline} we only \textbf{read} \code{ser},
but do \textbf{not modify} it. Python allows access to global variables
\textbf{without a} \code{global}\textbf{ declaration, as long as you only
read them}. Only when you try to modify them (as with
\code{buffer += ...}) do you need the \code{global}.\\
This is a nasty but important distinction!
\end{region}

\begin{region}[pitch=13.3]{648.6}
\begin{ptasks}[start=5]
\item Even though it is \textbf{not technically necessary}, change
  \textbf{line 10} as follows:
\end{ptasks}
\end{region}

\begin{region}[x=74.7,right=0]{661.9}
\begin{codeblock}[style=py,basicstyle=\codex{12}]
@10 @    @global@ buffer, ser
\end{codeblock}
\end{region}

\begin{region}[x=74.7,right=0,pitch=14.3]{676.0}
For ordinary mortals, it is then immediately \textbf{clear} that this
function accesses \textbf{global variables} -- whether for reading or
writing. This makes the code \textbf{easier to understand} and protects
against errors if you later do change something about \code{ser} after all.
\end{region}

\begin{region}[pitch=13.4]{717.9}
\begin{ptasks}[start=6]
\item Now add the new file to the repository:\\
  \code{git add controller.py}\\
  \code{git commit -a -m "proof of concept"}\\
  \code{git push}
\end{ptasks}
\end{region}

\nextpage
%% ===================================================================== 8 ==
\begin{region}{69.4}
\heading{Step 4: Extending the proof of concept}
\end{region}

\begin{region}[pitch=13.75]{85.2}
We can now copy code from our program \code{serial-reader.py} into
\code{controller.py} to obtain the \textbf{actual functionality}. This
involves a bit of typing -- but it is not particularly tricky.
\end{region}

\begin{region}[pitch=14.3]{129.0}
Below you can see what the code might look like \textbf{after a first
iteration}. First, a brief overview of \textbf{what has been changed or
added}:
\end{region}

\begin{region}[pitch=13.8]{158.8}
\begin{dashlist}
\item \textbf{Lines 1--2:} The required modules are imported:
  \begin{dashlist}
  \item \code{serial}, to communicate with the serial interface,
  \item \code{time}, to put the process to sleep from time to time.
  \end{dashlist}
\item \textbf{Lines 7--11:} The function \code{connect} now actually
  establishes a connection to the serial interface and initializes
  \code{ser}.
\item \textbf{Lines 13--27:} The function \code{readline} reads all
  available characters, appends them to \code{buffer}, and checks for the
  end of a line:
  \begin{dashlist}
  \item If a \textbf{line feed} (\code{\n}) is found, it prints the content
    of \code{buffer} and resets \code{buffer}.
  \item If a \textbf{carriage return} (\code{\r}) is found, nothing happens.
  \item Otherwise, the character is added to \code{buffer}.
  \item Important: The function \textbf{always} returns \code{None} -- even
    if it succeeds! (We will fix that later.)
  \end{dashlist}
\item \textbf{Lines 28--33:} The main program now uses \code{connect} and
  \code{readline}:
  \begin{dashlist}
  \item In \textbf{line 28}, the connection is established.
  \item Then \code{readline} is called in an infinite loop.
  \item If \code{readline} returns a line (that is, not \code{None}), it is
    output.
  \end{dashlist}
\end{dashlist}
\end{region}

\begin{region}[pitch=13.4]{395.8}
This already looks quite usable -- but: it does not quite work yet. Because,
as mentioned above, \code{readline} always returns \code{None}, even if a
line has actually been read.
\end{region}

\begin{region}[pitch=11]{424.1}
\begin{codeblock}[style=py,basicstyle=\codex{11}]
@  1 @!import! serial
@  2 @!import! time
@  3@
@  4 @ser = ~None~
@  5 @buffer = ~None~
@  6@
@  7 def@ ~connect~(port):
@  8 @    @global@ ser, buffer
@  9 @    ser = serial.Serial(port, ^9600^)
@ 10 @    buffer = ^""^
@ 11 @    ~print~(f^"connected to {port}"^)
@ 12@
@ 13 def@ ~readline~():
@ 14 @    @global@ ser, buffer
@ 15@
@ 16 @    @while@ ser.in_waiting > ^0^:
@ 17 @        ch = ser.read(^1^)
@ 18 @        ch = ~bytes~.decode(ch)
@ 19 @        @if@ ch == ^'^!\n!^'^:
@ 20 @            ~print~(f^"Zeile> {buffer}"^)
@ 21 @            buffer = ^""^
@ 22 @        @elif@ ch == ^'^!\r!^'^:
@ 23 @            @pass@
@ 24 @        @else@:
@ 25 @            buffer += ch
@ 26 @    @return@ ~None~
@ 27@
@ 28 @connect(^"/dev/ttyUSB0"^)
@ 29 while@ ^1^:
@ 30 @    line = readline()
@ 31 @    @if@ line @is@ @not@ ~None~:
@ 32 @        ~print~(f^"readline {line}"^)
@ 33 @    time.sleep(^0.1^)
\end{codeblock}
\end{region}

\nextpage
%% ===================================================================== 9 ==
\begin{region}{69.4}
\subheading{Tasks}
\end{region}

\begin{region}[pitch=14.25]{96.2}
\begin{ptasks}
\item Adapt \code{controller.py} accordingly. Confirm that the received data
  is not output by the \code{print} in \textbf{line 32} (as actually
  intended), but by the output in \textbf{line 20}.
\end{ptasks}
\end{region}

\begin{region}[pitch=13.45]{138.9}
\begin{ptasks}[start=2]
\item The obvious idea is to return the content of \code{buffer} in
  \textbf{line 20} instead of just outputting it. So change the line to
  \code{return buffer}. Confirm that the behavior has now changed, but is
  not yet quite what we want.
\end{ptasks}
\end{region}

\begin{region}[pitch=13.5]{179.9}
\begin{ptasks}[start=3]
\item \textbf{Idea for solving the current problem}\\
  When we detect a complete line in \code{readline()} (that is, at the
  \code{\n}), then we:
  \begin{Romans}
  \item store the current value of \code{buffer} in a new local variable,
    e.g., \code{line}
  \item reset \code{buffer} to the empty string \code{""}
  \item then return \code{line}
  \end{Romans}
\end{ptasks}
\end{region}

\begin{region}[x=74.7,right=0,pitch=13.4]{261.4}
\textbf{Think for yourselves} about how you can implement this strategy in
Python code. Adapt the corresponding part of the function
\code{readline()} accordingly. Then run the program and check whether
everything works as desired.
\end{region}

\begin{region}[pitch=13.55]{315.5}
\begin{ptasks}[start=4]
\item \textbf{If everything now works as desired:}\\
  \code{git commit -a -m "connect, readline"}\\
  \code{git push}
\end{ptasks}
\end{region}

\begin{region}{369.9}
\heading{Step 5: More structure -- er, make it classy!}
\end{region}

\begin{region}[pitch=13.9]{391.9}
In our program, certain things \textbf{obviously belong together}:
\begin{dashlist}
\item the two global variables \code{ser} and \code{buffer}, and
\item the functions \code{connect} and \code{readline}.
\end{dashlist}
We \textbf{know} this -- at least \textbf{today}, because we worked hard to
get these four loners to finally work together \textbf{as a team}.
\end{region}

\begin{region}[pitch=14.2]{474.9}
But: if you do not yet know the code and how it works, you do \textbf{not
immediately} see that these things belong together.
\end{region}

\begin{region}[pitch=13.6]{516.6}
Fortunately, Python offers a language feature that lets you \textbf{make}
exactly this structure \textbf{visible}. And yes -- this is where the magic
word comes up: \code{class}
\end{region}

\begin{region}[pitch=13.8]{553.1}
Here is an almost self-explanatory example of how our code can be expressed
\textbf{more clearly and in a more structured way} with a class. And yes --
it is once again a \textbf{proof of concept}, which you will transfer to the
real code in a moment:
\end{region}

\begin{region}[pitch=12]{607.1}
\begin{codeblock}[style=py,basicstyle=\codeix]
@  1 class@ ~Controller~:
@  2 @    ser = ~None~
@  3 @    buffer = ~None~
@  4@
@  5 @    @def@ ~connect~(port):
@  6 @        Controller.ser = f^"connected to {port}"^
@  7 @        Controller.buffer = ^"empty"^
@  8@
@  9 @    @def@ ~readline~():
@ 10 @        ~print~(f^"reading from {Controller.ser}"^)
@ 11 @        Controller.buffer += ^", some data"^
@ 12 @        ~print~(f^"buffer: {Controller.buffer}"^)
@ 13@
@ 14 @Controller.connect(^"bar"^)
@ 15 @Controller.readline()
\end{codeblock}
\end{region}

\nextpage
%% ==================================================================== 10 ==
\begin{region}{69.4}
\subheading{Tasks}
\end{region}

\begin{region}[pitch=13.2]{97.3}
\begin{ptasks}
\item Write the example program above into a file named
  \code{controller_test.py} and run the program.
\end{ptasks}
\end{region}

\begin{region}{123.7}
\begin{ptasks}[start=2]
\item You should have received the following output:
\end{ptasks}
\end{region}

\begin{region}[x=74.7,right=0,pitch=11]{148.7}
\begin{codeblock}[basicstyle=\fontsize{9bp}{11bp}\ttfamily]
reading from connected to bar
buffer: empty, some data
\end{codeblock}
\end{region}

\begin{region}[x=74.7,right=0]{186.2}
And with that, you can \textbf{already guess} what this small program is
meant to express:
\end{region}

\begin{region}[x=74.7,right=0,pitch=13.5]{212.9}
The variables \code{ser} and \code{buffer} as well as the functions
\code{connect} and \code{readline} are part of the class \code{Controller}.
\end{region}

\begin{region}[x=74.7,right=0,pitch=13.75]{252.8}
One could also say:
\begin{dotlist}
\item They are encapsulated in the class,
\item or they belong to the namespace \code{Controller}.
\end{dotlist}
\end{region}

\begin{region}[x=74.7,right=0,pitch=13.6]{305.9}
The world of object-oriented programming (OOP) may be full of clever
terms -- but what is meant is basically quite simple:
\begin{dotlist}
\item The variable \code{ser} in line 2 is actually called:
  \code{Controller.ser}
\item The variable \code{buffer} in line 3 is called: \code{Controller.buffer}
\item The function \code{connect} in line 5 is fully called:
  \code{Controller.connect}
\item And \code{readline} in line 9 is: \code{Controller.readline}
\end{dotlist}
\end{region}

\begin{region}[x=74.7,right=0]{403.8}
\emoji{1F449} \textbf{And that is exactly why we no longer need a}
\code{global}\textbf{!}
\end{region}

\begin{region}[x=74.7,right=0,pitch=13.8]{420.5}
Because with this language feature (\code{class}), it is \textbf{absolutely
clear} \textbf{which variable} the code refers to -- namely the one that
\textbf{belongs to the class} \code{Controller}. This provides clarity,
structure -- and fewer surprises.
\end{region}

\begin{region}[x=74.7,right=0,pitch=17.6]{474.7}
So -- now it should be easy to apply this principle to \code{controller.py}.
And once again, it is your turn! \emoji{1F60E}
\end{region}

\begin{region}[pitch=13.5]{508.0}
\begin{ptasks}[start=3]
\item If everything works as before (only now with class):\\
  \code{git commit -a -m "mit class"}\\
  \code{git push}
\end{ptasks}
\end{region}

\end{document}
